\section{Strategy Comparison at Scale}
\label{sec:strategy}

The calibration (\S\ref{sec:calibration}) and single-pass experiments
(\S\ref{sec:singlepass}) established \textit{how} recall
works and \textit{what} compaction destroys. This section addresses the core question:
when compaction must be applied
repeatedly over a long conversation, \textbf{which strategy preserves the most
information?}

\subsection{Strategies under evaluation}
\label{sec:strategies}

\textbf{Brutal (single-shot)}: When context exceeds 90\% of max, summarize ALL
messages except the 2 most recent. The input is truncated to a character cap
($\sim$150K real tokens) before summarization.

\begin{figure}[h]
\centering
\includegraphics[width=0.85\linewidth]{fig_strategy_brutal.png}
\caption{Brutal: a single summarization call replaces almost the entire context with one summary plus the last two messages.}
\label{fig:strategy-brutal}
\end{figure}

\textbf{Incremental (dual watermark)}: When context exceeds 90\%, compact enough old
messages to bring context down to 60\%. Previous summaries ARE included in the
text to be re-summarized. Creates a ``JPEG cascade'' where each cycle degrades
earlier summaries.

\begin{figure}[h]
\centering
\includegraphics[width=0.85\linewidth]{fig_strategy_incremental.png}
\caption{Incremental: every cycle re-summarizes the prior summary together with the oldest raw messages, accumulating loss across iterations (the ``JPEG cascade'').}
\label{fig:strategy-incremental}
\end{figure}

\textbf{Frozen (dual watermark + immutable summaries)}: Same trigger and target as
incremental, but completed summaries are marked as frozen and never
re-summarized. Only raw (non-frozen) messages are compacted. When frozen
summaries exceed a budget, the oldest are merged.

\begin{figure}[h]
\centering
\includegraphics[width=0.85\linewidth]{fig_strategy_frozen.png}
\caption{Frozen: only raw messages are compacted; completed summaries are immutable, accumulating in a budgeted ``frozen'' zone at the bottom of the context.}
\label{fig:strategy-frozen}
\end{figure}

\textbf{FrozenRanked (hierarchical merge)}: A variant of Frozen where each frozen
summary carries a \textit{rank} (initially 1). When frozen summaries exceed their
budget, only two summaries of the \textit{lowest available rank} merge --- producing
a summary of rank+1. This limits each fact to at most $\log_2(N)$ compression
passes, compared to $N/2$ in Frozen's sequential oldest-first merging.

\begin{figure}[h]
\centering
\includegraphics[width=0.85\linewidth]{fig_strategy_frozenranked.png}
\caption{FrozenRanked vs Frozen: instead of always merging the oldest pair (left), FrozenRanked merges the two lowest-rank summaries (right), bounding the number of times any single summary is re-compressed to $\log_2 N$ rather than $N/2$.}
\label{fig:strategy-frozenranked}
\end{figure}

The key insight: in Frozen, the oldest summary accumulates \textit{all} merges
sequentially (cascade). In FrozenRanked, merges are balanced like a
tournament bracket --- summaries of similar compression depth merge together,
preventing any single summary from being re-compressed excessively.

\subsection{Protocol: a constant-density, replicated comparison}
\label{sec:strategy-protocol}

Pilot studies (synthetic facts at 1.5M and 3M tokens, then a first
LongMemEval-grounded study with 80 fixed facts spanning 500K--10M) suggested
the ordering Frozen $>$ Incremental $>$ Brutal but two design issues prevented
clean conclusions. First, as the conversation grew, fact density $\delta$
decreased mechanically (0.16 $\to$ 0.008 facts/kTok), so degradation reflected
both increased compaction cycles \textit{and} fact dilution. Second, each cell was
measured once: with a non-deterministic compaction LLM (see \S\ref{sec:methodological} below),
single measurements proved insufficient to separate signal from noise.

The experiment described here resolves both issues. A single 5M conversation
is built with \textbf{constant fact density} (200 facts at $\delta=0.04$ facts/kTok,
uniformly interleaved between 0.5\% and 99.6\% of the tokens). The same
conversation is fed to all four strategies and evaluated at five mid-feed
checkpoints (500K, 1M, 2M, 3.5M, 5M) so that older facts age under repeated
compaction while newer facts are still fresh. Each (strategy, checkpoint)
cell is measured \textbf{4 to 6 times} across independent runs.

Note that $\delta=0.04$ is well below the calibration ``sweet spot''
($\sim$0.42, \S\ref{sec:calibration}) at
which static-context recall plateaus. With density this low, absolute recall
levels in \S\ref{sec:strategy} are bounded by limited evidence as well as by compaction loss.
The strategy \textit{ordering} (S4 $>$ S3 $>$ S2 $>$ S1) is the salient comparison; the
absolute numbers are not directly comparable to \S\ref{sec:calibration}/\S\ref{sec:singlepass} baselines.

\textbf{Setup summary}:

\begin{table}[h]
\centering
\caption{Strategy comparison setup at $\delta=0.04$.}
\label{tab:strategy-setup}
\begin{tabular}{ll}
\toprule
Parameter & Value \\
\midrule
Conversation & 5M tokens (4.9M actual), built once with seed 42 \\
Facts        & N=200 from R4, uniformly interleaved (0.5\% $\to$ 99.6\%) \\
Density      & $\delta = 0.04$ facts/kTok, constant \\
Checkpoints  & 500K, 1M, 2M, 3.5M, 5M (mid-feed) \\
Strategies   & S1, S2, S3, S4 (same input conversation) \\
QA model     & Claude Sonnet 4.6, $Q=5$ questions per prompt \\
Judge model  & Claude Haiku 4.5, strict prompt (see \S\ref{sec:methodological}) \\
Sampling     & API temperature parameter set to 0 \\
Replicates   & 4--6 independent runs per (strategy, checkpoint) cell \\
\bottomrule
\end{tabular}
\end{table}

The replicates are spread across batch-API runs and gateway runs (some
gateway runs covered S3+S4 only, others S1+S2 only). Some cells additionally
include QA-only replicates: the saved compacted contexts are re-evaluated
with a fresh Q\&A pass and a new judge call, isolating the variance
contributed by the QA / judge phase from the variance contributed by
compaction.

API ``temperature 0'' deserves a note here. Setting \code{temperature=0} is the
conventional way to ask for a deterministic-looking output from these
models; in practice it makes the model strongly favour a single high-
probability token at each step. Across \textbf{independent sessions} with the
same input it does \textit{not} yield identical outputs: small differences
(implementation, infrastructure load, internal session state) can change
the chosen token. Our compaction calls all use temperature 0 yet produce
visibly different summaries on identical inputs --- see \S\ref{sec:methodological} for the
quantitative impact.

\subsection{Results}
\label{sec:strategy-results}

\begin{figure}[h]
\centering
\includegraphics[width=0.85\linewidth]{fig_phase_d_final.png}
\caption{Recall (mean $\pm$ std across replicates) for each strategy as a
function of conversation length, evaluated on a single 5M conversation at
five mid-feed checkpoints. n=4--6 replicates per point. $\delta=0.04$ throughout.
Sonnet 4.6 QA + Haiku 4.5 strict judge.}
\label{fig:phase-d-final}
\end{figure}

The mean recall and standard deviation across replicates are
(QA: Claude Sonnet 4.6, $Q=5$; Judge: Claude Haiku 4.5, strict prompt):

\begin{table}[h]
\centering
\caption{Mean $\pm$ std recall across replicates for each strategy and checkpoint.}
\label{tab:phase-d-mean}
\small
\begin{tabular}{cccccc}
\toprule
Checkpoint & S1 Brutal & S2 Incremental & S3 Frozen & S4 FrozenRanked \\
\midrule
500K (n=6)  & 7.0 $\pm$ 2.5\% & 16.7 $\pm$ 6.4\% & 14.9 $\pm$ 7.7\%  & \textbf{28.1 $\pm$ 12.0\%} \\
1M (n=6)    & 6.0 $\pm$ 1.9\% & 9.4 $\pm$ 5.7\%  & 13.2 $\pm$ 15.0\% & \textbf{13.7 $\pm$ 14.0\%} \\
2M (n$\approx$6) & 0.5 $\pm$ 0.6\% & 0.8 $\pm$ 1.5\%  & 4.6 $\pm$ 5.3\%   & \textbf{11.2 $\pm$ 5.9\%} \\
3.5M (n=5)  & 4.0 $\pm$ 2.2\% & 6.8 $\pm$ 4.0\%  & 6.1 $\pm$ 5.7\%   & \textbf{9.4 $\pm$ 6.3\%} \\
5M (n=4)    & 2.5 $\pm$ 0.0\% & 2.8 $\pm$ 0.5\%  & 3.0 $\pm$ 2.9\%   & \textbf{4.6 $\pm$ 3.6\%} \\
\bottomrule
\end{tabular}
\end{table}

Four patterns emerge from this dataset:

\textbf{1. The strategy hierarchy holds in the means.} S4 FrozenRanked has the
highest mean recall at every checkpoint, followed by S3 Frozen, S2
Incremental, and S1 Brutal. To distinguish robust orderings from
within-noise differences we ran Wilcoxon signed-rank tests on paired
replicates (one-sided, alternative S4 $>$ S3 etc.); the table below reports
the n paired replicates available per cell, the mean difference, and the
one-sided Wilcoxon p-value. Sample sizes are limited (n=2--6) and
1M/3.5M variance is large, so several gaps fall short of significance
even where the mean ordering is consistent.

\begin{table}[h]
\centering
\caption{Pairwise Wilcoxon signed-rank tests for strategy ordering.}
\label{tab:strategy-wilcoxon}
\small
\begin{tabular}{lccccc}
\toprule
Comparison & 500K & 1M & 2M & 3.5M & 5M \\
\midrule
S4 vs S3 & +13.2pp, \textbf{p=0.03} (n=6) & +0.4pp, n.s.\ (n=6) & +6.7pp, \textbf{p=0.02} (n=6) & +2.0pp, n.s.\ (n=5) & +1.6pp, p=0.06 (n=4) \\
S4 vs S1 & +13.2pp, p=0.06 (n=4) & +6.4pp, n.s.\ (n=4) & +9.7pp, p=0.06 (n=4) & +5.5pp, p=0.06 (n=4) & 0.0pp, n.s.\ (n=2) \\
S2 vs S1 & +9.7pp, \textbf{p=0.03} (n=6) & +3.4pp, n.s.\ (n=6) & +0.3pp, n.s.\ (n=5) & +3.4pp, n.s.\ (n=4) & +0.3pp, n.s.\ (n=3) \\
\bottomrule
\end{tabular}
\end{table}

(Full pairwise table in \code{phase\_d\_stats.json}.) \textbf{In summary}: the mean
ordering S4 $>$ S3 $\geq$ S2 $>$ S1 is consistent across all checkpoints, but
robustly significant only at the lowest and middle scales (S4 $>$ S3 at 500K
and 2M with p$<$0.05, all S4 $>$ Sk wins-only outcomes at 2M). At 1M
and at 5M the standard deviations are too large for n=4--6 replicates to
yield significance.

\textbf{2. The S4--S3 gap is real at small/intermediate scales but narrows at 5M.}
Pilot experiments (\S\ref{sec:strategy-protocol}) and the older fixed-density study showed
S3 $\approx$ S4 at all sizes because frozen summaries never accumulated enough to
trigger merges. With constant density and replicates, S4 has a robust
advantage at 500K (28.1 vs 14.9, +13.2pp, \textbf{p=0.03}), 2M (11.2 vs 4.6,
+6.7pp, \textbf{p=0.02}), and a marginal one at 5M (4.6 vs 3.0, +1.6pp,
p=0.06; 4/4 wins). At 1M the gap effectively disappears (+0.4pp, n.s.),
which we attribute to one outlier replicate where S3 happened to recall
unusually well. We interpret the overall pattern as FrozenRanked balancing
the merge load across the frozen chain (cf.\ \S\ref{sec:strategies}), preserving recall
while the chain is short enough for the bracket structure to matter;
the persistent (if marginal) S4--S3 gap at 5M suggests this is not purely
a small-scale phenomenon.

\textbf{3. A tentative non-monotonic dip around 2M.} Mean recall for S1 and S2
is markedly lower at 2M (S1=0.5\%, S2=0.8\%) than at the surrounding 1M
(6.0\%, 9.4\%) and 3.5M (4.0\%, 6.8\%) checkpoints; S3 (4.6\%) and S4 (11.2\%)
also bottom out at 2M. One plausible interpretation is a \textbf{transient
overcompression effect}: at the 2M checkpoint, the most recent compaction
cycle has just consumed a large block of older facts whose summaries have
not yet stabilised, and few fresh facts have been fed in since. By 3.5M,
new facts populate the recent window and the recall partially recovers.
However, the standard deviations at 2M overlap with neighbouring
checkpoints (e.g.\ S2 $\sigma=1.5$pp at 2M vs $\sigma=4.0$pp at 3.5M), and the effect
rests on a single conversation; we therefore report this as a \textbf{tentative
observation} that needs replication on additional conversation seeds
(see \S\ref{sec:future}).

\textbf{4. All strategies converge near 5\%--10\% recall at 5M.} Even the best
strategy (S4) retains only 4.6\% of the 200 facts. This convergence
mirrors \S\ref{sec:singlepass}'s single-pass result: at high cumulative compaction (50--80\%
of the original content compacted), the differences between strategies
become small relative to the absolute amount of information lost.
The gap between S1 (2.5\%) and S4 (4.6\%) at 5M is 2.1pp --- non-zero in
mean but well within one standard deviation.

\textbf{Spatial recall: erosion of older facts.} To complement the global recall
trajectory, we trace recall as a function of fact position in the original
conversation, with one curve per checkpoint per strategy:

\begin{figure}[h]
\centering
\includegraphics[width=0.95\linewidth]{fig_phase_d_spatial_per_strategy.png}
\caption{Spatial recall per strategy. Each panel is one strategy; each
curve is one checkpoint (averaged across all replicates). The x-axis is the
position of the fact in the full 5M-token conversation. Earlier checkpoints
only cover the first part of the conversation; later checkpoints extend
further right. The decay of older positions across successive checkpoints
visualizes erosion under repeated compaction.}
\label{fig:phase-d-spatial}
\end{figure}

For S1 Brutal and S2 Incremental, recall is concentrated near the \textbf{right
edge} of each curve --- the most recent facts at the moment of evaluation.
Earlier positions, well-sampled at earlier checkpoints, drop to near zero
by the time the conversation reaches 3.5M or 5M. This is the JPEG cascade
visualised: each compaction cycle erases a slice of older content.

For S3 Frozen and S4 FrozenRanked, the curves retain non-zero recall in
the \textbf{left half} of each panel even at later checkpoints --- frozen
summaries continue to make early facts retrievable, though at a much
reduced level than when those facts were fresh. The S4 panel is visibly
flatter than the S3 one across the conversation, which is the spatial
counterpart to its higher mean recall.

\subsection{Methodological findings: variance and judge sensitivity}
\label{sec:methodological}

Two methodological observations from this study are worth highlighting
because they affect how results from any compaction benchmark should be
interpreted.

\textbf{Compaction is not deterministic.} Even at temperature 0, Sonnet produces
different summaries when the same compaction call is invoked in different
sessions. Across our independent runs on identical conversations and seeds,
the resulting frozen summaries differ in length and content (e.g., the same
S4 checkpoint has 639 messages but 658K vs 630K characters across two
runs), and the downstream recall measurements differ accordingly. For S4
at 1M the recall ranges from 2.6\% to 35.9\% across replicates --- a factor of
14$\times$. The mean across replicates is stable, but a single measurement of
``strategy A vs strategy B'' can be misleading by an order of magnitude in
either direction. \textbf{Multiple replicates per cell are required} to make
reliable claims about strategy ordering.

\textbf{Judge prompt sensitivity matters.} Our initial judge prompt (Haiku 4.5,
``recalled = does the answer demonstrate knowledge of the expected
information?'') proved too lenient: when the LLM answer was \textit{``I don't recall
discussing X, but Y was mentioned''}, Haiku frequently scored
\code{recalled = true} because the answer mentioned a related topic. Re-judging
with a stricter prompt (recalled = true only if the \textit{specific expected
fact} is present in the answer; ``I don't recall'' or generic mentions of
related topics are always false) reduced recall numbers by 5--15pp at
intermediate checkpoints. All numbers reported above use the strict judge.
The standard prompt remains in the codebase for reproduction of the
\S\ref{sec:calibration}--\S\ref{sec:singlepass}
calibration results, which used the same lenient judge but where the
looseness has less effect because answers are either clearly correct or
clearly absent (less ambiguity).

\textbf{Human-judge calibration.} To validate the strict judge, one of the
authors hand-graded two stratified samples of 50 (question, expected,
LLM-answer) tuples each:

\begin{itemize}
  \item \textbf{Sample A --- random across all answers}: dominated by ``I don't recall''
    items (the bulk of the corpus). Validates the \textit{negative} class.
  \item \textbf{Sample B --- items where lenient OR strict said \code{recalled = True}}:
    validates the \textit{positive} class.
\end{itemize}

Cohen's $\kappa$ between the human and each LLM judge:

\begin{table}[h]
\centering
\caption{Cohen's $\kappa$ between human grader and the two LLM judges.}
\label{tab:human-judge-kappa}
\small
\begin{tabular}{lcc}
\toprule
 & Sample A (n=50, mostly negatives) & Sample B (n=50, positives) \\
\midrule
\textbf{$\kappa$ Human $\leftrightarrow$ Strict} (recalled) & \textbf{1.000} & \textbf{0.831} \\
$\kappa$ Human $\leftrightarrow$ Lenient (recalled) & \makecell[c]{0.000\\(2 lenient FP)} & \makecell[c]{0.000\\(lenient says True for all 50\\by construction $\to$ no variance)} \\
$\kappa$ Human $\leftrightarrow$ Strict (accurate) & 1.000 & 0.879 \\
$\kappa$ Human $\leftrightarrow$ Lenient (accurate) & 1.000 & 0.716 \\
\bottomrule
\end{tabular}
\end{table}

The strict judge is in \textbf{almost-perfect agreement} with the human
($\kappa \geq 0.83$ on both samples; Landis \& Koch threshold: $\kappa \geq 0.81$). The
disagreements between human and strict on Sample B are interpretable:
they are mostly cases where the LLM says \textit{``I don't recall, but from the
summary your sales were \dots''} and the strict judge accepts the bracketed
hint while the human treats only the explicit ``I don't recall'' as
not-recalled. There are also 1--2 cases where the LLM provides a wrong
specific value (e.g.\ expected ``120 stars'', answered ``500 stars''): the
strict judge marks \code{recalled = True, accurate = False} per its rule
(specific number mentioned, wrong value), while the human prefers
\code{recalled = False} when the value is wrong. This is a genuine definitional
ambiguity (recalled = ``topic-correct'' vs ``value-correct''), not a
mis-calibration.

In contrast, the lenient judge is essentially uninformative on the
positive class: it labels 100\,\% of Sample B as \code{recalled = True} (by
construction --- Sample B is filtered on \code{lenient\_recalled = True}), so
the kappa with the human collapses to zero. On Sample A the lenient
judge produces 2 visible false positives where the LLM says ``I don't
recall mentioning X'' --- the very pattern the strict prompt was designed
to catch.

The samples and per-item judgements are in the repository at
\code{human\_judge\_sample.json} (Sample A) and \code{human\_judge\_positives.json}
(Sample B). All \S\ref{sec:strategy} numbers in this paper use the strict judge.

\textbf{Spatial / quintile breakdowns.} With 200 facts, mid-feed checkpoints,
and replicates, per-quintile breakdowns become noisy at small recall
levels. The structural trade-off identified in the pilot --- Frozen
preserves the past at the cost of the present, Incremental does the
opposite --- qualitatively persists in this study, but the small absolute
counts (often 0--3 recalled facts per quintile) make per-cell tables
less informative than the integrated mean $\pm$ std presented above.

\subsection{A second-seed generalization attempt --- and why it failed}
\label{sec:seed99}

Independent peer review of an earlier draft correctly flagged that all \S\ref{sec:strategy}
results come from a \textit{single} 5M conversation built with seed 42. To
address the generalization concern, we built a second 5M conversation
with seed 99 (same density, same fact bank, different LongMemEval
sessions used as padding) and re-ran the four strategies through the
same checkpoint pipeline.

The run completed but the recall numbers are uninterpretable as a
generalization test. At checkpoint 502K, S3 and S4 had executed \textbf{0
compaction cycles} (the trigger threshold was not yet reached) and
fed all $\sim$2{,}180 raw messages to the QA model. The QA model used
(\code{claude-sonnet-4-20250514}, 200K context window) silently
truncated/rejected these contexts, producing batch errors and a
flat 0\% recall for S3/S4 at every checkpoint. S1, which compacts
aggressively from the start, fits inside the QA window and shows
recall in the 1--5\% range --- but with so few successful answers the
ordering is meaningless.

Two lessons from this failure are worth recording.

\textbf{The instrumentation has a budget envelope that is easy to break.}
The \S\ref{sec:strategy} measurements implicitly assume that the compacted-context size is $\leq$ the QA
model window. With seed 42 this held by construction; with seed 99 the
mix of LongMemEval sessions produced longer messages on average, so
S3/S4 (which delay compaction the most) crossed the QA model's window
before the compaction trigger fired. Future replications must either
(a) lower the strategies' compaction trigger so the post-compaction
context fits the QA window with margin, or (b) use a QA model with a
1M-token window (Sonnet 4.5 / 4.6) --- at materially higher cost.

\textbf{Single-seed \S\ref{sec:strategy} results are an honest admission, not a hidden weakness.}
The methodological consequence is that the strategy hierarchy reported
in \S\ref{sec:strategy-results} (S4 $>$ S3 $>$ S2 $>$ S1) is, strictly speaking, established on one
conversation, with the variance numbers reflecting compaction-LLM
stochasticity but not conversation-level generalization. We document
the seed-99 run in the repository as \code{iterative\_v6\_R4\_20260503\_1638/}
with this caveat: it is a budget-and-instrument failure to be fixed
before generalization claims can be made, not a counter-finding. The
pilot studies on synthetic facts (1.5M and 3M tokens, briefly
described in \S\ref{sec:strategy-protocol}) showed the same qualitative S4 $>$ S3 $>$ S2 $>$ S1
ordering, which is the strongest cross-conversation evidence currently
available; a properly instrumented multi-seed study is left to future
work (\S\ref{sec:future}).
